\documentclass[10pt,a4paper]{amsart}
\usepackage[margin=19mm]{geometry}
\usepackage{amsmath,amssymb,mathtools}
\usepackage[scr=rsfs,cal=euler]{mathalfa}
\usepackage{xfrac}
\usepackage[unicode,hidelinks,bookmarks=false]{hyperref}
\newcommand{\ie}{i.e.\ }
\newcommand{\cf}{cf.\ }
\newcommand{\resp}{respectively\ }
\newcommand{\Subsection}{\subsection}
\providecommand{\nobreakdash}{\nobreak\hbox{-}}
\setlength{\emergencystretch}{2em}
\multlinegap0pt
\usepackage{amssymb}
\newtheorem{lem}{Lemma}
\def \N {\mathbb N}
\def\a{\alpha}
\def\b{\kern 0.02em \flat \kern 0.02em}
\def\m{\mathcal M}
\title{Typical periodic optimization for dynamical systems: Symbolic dynamics}
\author{Wen Huang, Oliver Jenkinson, Leiye Xu, Yiwei Zhang}
\date{Source: arXiv:2603.07224v1 (CC BY 4.0)}
\begin{document}
\maketitle

\noindent\emph{Source and licence.} Typical periodic optimization for dynamical systems: Symbolic dynamics. Version arXiv:2603.07224v1 (CC BY 4.0), distributed by arXiv under the Creative Commons Attribution 4.0 International licence, http://creativecommons.org/licenses/by/4.0/, as recorded in the arXiv licence metadata for this exact version and shown on the abstract page https://arxiv.org/abs/2603.07224v1. Full licence notice: \texttt{licenses/arxiv-2603.07224-CC-BY-4.0.txt}.\par\smallskip

\noindent\textit{Editorial scope.} Counted unit: the article's lem 2\_to\_the\_n (source0.tex lines 3900--3910). The article's complete original proof of it is delivered with it (source0.tex lines 3911--4023, one \texttt{proof} environment of 113 lines). No mathematics is written, repaired or re-derived: the statement and the proof are the article's own lines, in the article's own notation, and no retained line is substituted or re-flowed.  Retained uncounted as the source-local context the proof needs: lines 3877--3890.  Nothing was omitted from any retained span, so the package carries no omission record.  No retained line contains a citation, so the wrapper carries no bibliography and no bibliography entry.  No retained line contains a figure environment, an image-inclusion command or drawing code, so no image is delivered and the package declares no asset dependency.  Editorial formatting only, in addition: an AMS \texttt{amsart} preamble that restates the article's own declarations and macros used by the retained lines, namely the environments \texttt{lem} on separate counters, together with the standard packages those lines need.  The retained environments print in this document as Lemma 1 in order of appearance, whereas the article prints the same items with the numbering of its own full document.  The complete native source of the article as distributed in the arXiv e-print for this exact version is bundled unchanged as \texttt{sources/195926/source0.tex}.
\begin{lem}\label{wasserstein_lemma}
If $(Y,d)$ is a compact metric space, the Wasserstein distance $W$ satisfies
\begin{itemize}
\item[(a)] If $\mu,\nu\in\m(Y)$ then $W(\mu,\nu)\le \text{diam}(Y)$.
\item[(b)] If $x,y\in Y$ then $W(\delta_x,\delta_y)=d(x,y)$.
\item[(c)] If $n\in\N$, and  $a_i\ge0$ for $1\le i\le n$ with $\sum_{i=1}^n a_i=1$,  and $\mu_i,\nu_i\in\m(Y)$  for $1\le i\le n$,
then 
$$
W\left( \sum_{i=1}^n a_i \mu_i, \ \sum_{i=1}^n a_i \nu_i \right)
\le \
\sum_{i=1}^n a_i  W(\mu_i, \nu_i ).
$$
\end{itemize}
\end{lem}

\begin{lem}
If $\lambda$ is the Morse measure, and $x=\theta_n \b x'$ for some $n\in\N$, $x'\in X$, 
then
\begin{equation}\label{2_to_the_n}
W(E_{2^{n}}(x),\lambda)\le \left(\frac{2\a}{1-\a}\right) \frac{1}{2^{n}},
\end{equation}
and
\begin{equation}\label{2_to_the_n_plus_one}
W(E_{2^{n}+1}(x),\lambda)\le \left(\frac{1+\a}{1-\a}\right) \frac{1}{2^{n}+1}.
\end{equation}
\end{lem}

\begin{proof}
For each $n\in\N$, 
$\lambda$ is the weak$^*$ limit
$\lim_{M\to\infty} E_{2^nM}(\omega)$, and we can write
\begin{equation}\label{E_m_M}
E_{2^nM}(\omega)
=
\frac{1}{M}\sum_{m=0}^{M-1} E_{2^n}\left(\sigma^{m2^n}(\omega)\right),
\end{equation}
so that if it can be shown that
\begin{equation}\label{W_E_m}
W\left(E_{2^n}(x),  E_{2^n}\left(\sigma^{m2^n}(\omega)\right) \right) \le \left(\frac{2\a}{1-\a}\right) \frac{1}{2^{n}}
\end{equation}
for $0\le m\le M-1$, then (\ref{E_m_M}), (\ref{W_E_m}) and Lemma \ref{wasserstein_lemma} give
$$
W\left(E_{2^{n}}(x),  E_{2^nM}(\omega)\right)\le
\frac{1}{M} \sum_{m=0}^{M-1} W( E_{2^{n}}(x),  E_{2^n}\left(\sigma^{m2^n}(\omega)\right) )
\le \left(\frac{2\a}{1-\a}\right) \frac{1}{2^{n}},
$$
for all $M\in\N$, and hence (\ref{2_to_the_n}) follows from the fact that $\lambda=\lim_{M\to\infty} E_{2^nM}(\omega)$.

To prove (\ref{W_E_m}), each of the $2^n$ points in the support of $E_{2^n}(x)$ 
will be paired with a nearby point in the support of
$E_{2^n}(\sigma^{m2^n}(\omega))$, with the distance between these points 
equal to the Wasserstein distance between the corresponding Dirac measures, by Lemma \ref{wasserstein_lemma}(b).
Specifically, writing $\omega=\omega_1\omega_2\ldots$, so that
	$\omega=\theta^n(\omega)=\theta^n(\omega_1)\theta^n(\omega_2)\cdots$
for all $n\in\N$,
we see that 
 the length-$2^n$ prefix of
$\sigma^{m2^n}(\omega)$
is either $\theta^n(0)$ or $\theta^n(1)$,
for all $m\ge0$.

In the case that $\theta^n(0)=\theta_n$
is the length-$2^n$ prefix of
$\sigma^{m2^n}(\omega)$, since $\theta_n$ is also the length-$2^n$ prefix of $x=\theta_n \b x'$
then 
\begin{equation*}
d(\sigma^i(x),\sigma^i(\sigma^{m2^n}(\omega)))\le \alpha^{2^n-i} \quad\text{for }0\le i\le 2^n-1,
\end{equation*}
and therefore
$$
W\left(E_{2^n}(x),  E_{2^n}\left(\sigma^{m2^n}(\omega)\right) \right) \le
\frac{1}{2^n} \sum_{i=0}^{2^n-1} d(\sigma^i(x),\sigma^i(\sigma^{m2^n}(\omega)))
\le
\frac{1}{2^n} \sum_{i=0}^{2^n-1} \alpha^{2^n-i} < \frac{\alpha}{1-\alpha} \frac{1}{2^n},
$$
which in particular implies (\ref{W_E_m}).

In the case that $\theta^n(1)$
is the length-$2^n$ prefix of
$\sigma^{m2^n}(\omega)$,
since
$\theta^n(1)=\theta^{n-1}(1)\theta^{n-1}(0)$ and $\theta_n=\theta^{n-1}(0)\theta^{n-1}(1)$,
we see that $\sigma^{2^{n-1}}(\sigma^{m2^n}(\omega))$ and $x$
share the same length-$2^{n-1}$ prefix $\theta^{n-1}(0)=\theta_{n-1}$,
and
$\sigma^{m2^n}(\omega)$ and $\sigma^{2^{n-1}}(x)$
share the same length-$2^{n-1}$ prefix $\theta^{n-1}(1)$.
Consequently,
\begin{equation}\label{813}
d\left( \sigma^i(x), \sigma^i( \sigma^{2^{n-1}}(\sigma^{m2^n}(\omega)) )\right)
\le \alpha^{2^{n-1}-i} \quad\text{for }0\le i\le 2^{n-1}-1,
\end{equation}
and
\begin{equation}\label{814}
d\left( \sigma^i(\sigma^{2^{n-1}}(x)), \sigma^i(\sigma^{m2^n}(\omega) )\right)
\le \alpha^{2^{n-1}-i} \quad\text{for }0\le i\le 2^{n-1}-1,
\end{equation}
so the bounds (\ref{813}), (\ref{814}),
together with
an argument analogous to the one above, give
\begin{equation}\label{815}
W\left(E_{2^{n-1}}(x),E_{2^{n-1}}(\sigma^{2^{n-1}}(\sigma^{m2^n}(\omega)))\right) < \frac{\alpha}{1-\alpha}\frac{1}{2^{n-1}}
\end{equation}
and
\begin{equation}\label{816}
W\left(E_{2^{n-1}}(\sigma^{2^{n-1}}(x)),
E_{2^{n-1}}(\sigma^{m2^n}(\omega))\right) 
< \frac{\alpha}{1-\alpha}\frac{1}{2^{n-1}}.
\end{equation}
Now
\begin{equation*}
E_{2^n}(x)= \frac{1}{2}E_{2^{n-1}}(x) + \frac{1}{2}E_{2^{n-1}}(\sigma^{2^{n-1}}(x))
\end{equation*}
and
\begin{equation*}
E_{2^n}(\sigma^{m2^n}(\omega))
= \frac{1}{2} E_{2^{n-1}}(\sigma^{m2^n}(\omega))
+ \frac{1}{2} E_{2^{n-1}}(\sigma^{2^{n-1}}(\sigma^{m2^n}(\omega))),
\end{equation*}
so
\begin{align*}
&W \left(E_{2^n}(x),  E_{2^n}\left(\sigma^{m2^n}(\omega)\right) \right) \\
&\le
\frac{1}{2}W\left(E_{2^{n-1}}(x),E_{2^{n-1}}(\sigma^{2^{n-1}}(\sigma^{m2^n}(\omega)))\right) 
+
\frac{1}{2} W\left(E_{2^{n-1}}(\sigma^{2^{n-1}}(x)),
E_{2^{n-1}}(\sigma^{m2^n}(\omega))\right) ,
\end{align*}
and therefore (\ref{815}), (\ref{816}) yield the required inequality (\ref{W_E_m}).

Noting that $E_{2^n+1}(x)= \frac{2^n}{2^n+1} E_{2^n}(x)+\frac{1}{2^n+1} E_1(\b x')$, it follows from Lemma 
\ref{wasserstein_lemma}(c)
 that
$$
W(E_{2^{n}+1}(x),\lambda)\le  \frac{2^n}{2^n+1} W( E_{2^n}(x),\lambda) +\frac{1}{2^n+1} W(E_1(\b x'),\lambda),$$
and using (\ref{2_to_the_n}) together with the bound $ W(E_1(\b x'),\lambda)\le 1$
(see  Lemma \ref{wasserstein_lemma}(a))
 gives
(\ref{2_to_the_n_plus_one}).
\end{proof}

\end{document}
